% ==========================================================================
% Model refinements and their derivations
% ==========================================================================

\section{Model Refinements and Their Derivations}
\label{sec:refinements}

The reduced formulation of Section~\ref{sec:surrogates} was evaluated against a
patient-grounded three-dimensional reference (Section~\ref{sec:realdata}). Four
modifications to the reduced equations were required for that evaluation to be
well posed: a corrected growth law, an identifiable replacement for the
physics--residual weights, a state-dependent form of that replacement, and a
diagnostic that makes the resulting balance measurable. Each is derived below
together with the property it establishes, and
Table~\ref{tab:changes} summarises them.

\subsection{Growth law: from logistic to volumetric}
\label{sec:volumetric}

The reduced model of Eq.~\eqref{eq:ode_rt_2state} assumes logistic growth of the
tumour \emph{mass}. The reference dynamics, however, are the spatial integral of
a Fisher--Kolmogorov invasion, for which the growth law is different.

For a focal seed spreading through tissue, the reaction--diffusion front
advances at the constant Fisher--KPP speed
$v = 2\sqrt{Df}$, so the radius of the invaded region grows as
$R(t) = R_{0} + v\,t$. With the interior saturated at carrying capacity, the
integrated mass follows the invaded volume,
\begin{equation}
\label{eq:volumetric_mass}
y(t) \;=\; \int_{\Omega} u(x,t)\,dx \;\simeq\; c\,R(t)^{3}
\;=\; y_{0}\left(1 + b\,t\right)^{3},
\qquad b = \frac{v}{R_{0}} .
\end{equation}
Differentiating and eliminating $t$ through $R \propto y^{1/3}$ gives
\begin{equation}
\label{eq:volumetric_rate}
\dot y \;=\; 3\,c\,R^{2}\,\dot R \;=\; 3\,c^{1/3}v\;y^{2/3}
\;\;\propto\;\; y^{\,q}, \qquad q = \tfrac{2}{3},
\end{equation}
whereas the logistic law of Eq.~\eqref{eq:ode_rt_2state} is the case $q = 1$.
The exponent is therefore not a free modelling choice: it is fixed by the
geometry of a constant-speed invasion front, and $q=1$ over-predicts the
acceleration of an expanding tumour.

We accordingly generalise the tumour equation to
\begin{equation}
\label{eq:backbone_q}
\dot y \;=\; \rho\,y^{\,q}\!\left(1-\frac{y}{K_{\mathrm{eff}}}\right)
\;-\;\gamma\,z\,y\;-\;\mu\,y ,
\qquad
\dot z = -\frac{z}{\tau} + U(t),
\end{equation}
with $q$ either fixed at $2/3$ (\emph{volumetric} backbone), fixed at $1$
(the original \emph{logistic} backbone), or trained
(\emph{power} backbone, $q$ constrained to $[0.3,1.2]$ through a sigmoid).
Everything else in Eq.~\eqref{eq:ode_rt_2state} --- in particular the latent
damage state and the multi-fraction memory it provides --- is unchanged.

The mechanism is worth naming: for a sphere expanding at constant radial speed,
$\dot V = v\,|\partial V|\propto V^{2/3}$, so growth is proportional to the
\emph{surface} of the invaded region, not to its volume. The logistic law makes
it proportional to the volume, which is why it over-predicts the acceleration of
an expanding tumour.

\emph{What this achieves.} Measured on the untreated reference trajectories of
the real cohort, $y^{1/3}$ is linear in $t$ with median $R^{2} = 0.9995$
(against $0.9864$ for the mass itself), and the effective exponent recovered
from $\mathrm{d}\log\dot y / \mathrm{d}\log y$ has median $0.81$
($0.65$--$0.90$ across the cohort). The empirical value sits
between the idealised surface-growth exponent $2/3$ and the logistic $1$, and
clearly away from the latter (Fig.~\ref{fig:growth_law}); the offset from $2/3$
is expected, since in real anatomy the interior is not fully saturated and the
front is neither spherical nor of uniform speed. That the fixed value $2/3$
nevertheless outperforms a \emph{trained} exponent
(Table~\ref{tab:main_results}) is itself informative: the short assimilation
window does not determine $q$, so imposing the mechanistically derived value is
better than estimating it.

\subsection{Identifiability of the physics--residual weights}
\label{sec:gauge}

Eq.~\eqref{eq:pinode_weighted} introduces two weights, $\omega$ on the
mechanistic term and $s_{r}$ on the neural closure. Because they do not sum to
one they cannot be read as a convex mixture, and it is natural to ask what
balance they do express. The answer is that, individually, they express nothing:
both are exact gauge freedoms of the model.

\paragraph{The mechanistic weight.}
Expanded, the mechanistic tumour term is
\begin{equation}
\label{eq:frt_expanded}
f^{(y)}_{\mathrm{RT}}
= \rho\,y - \frac{\rho}{K_{\mathrm{eff}}}\,y^{2} - \gamma\,z\,y - \mu\,y ,
\end{equation}
which is \emph{jointly linear} in $(\rho,\gamma,\mu)$ at fixed
$K_{\mathrm{eff}}$ --- note that the quadratic coefficient $\rho/K_{\mathrm{eff}}$
carries the same $\rho$ as the linear one. It is therefore homogeneous of degree
one in those three parameters, so
\begin{equation}
\label{eq:omega_gauge}
\omega f^{(y)}_{\mathrm{RT}}(y,z,U;\rho,\gamma,K_{\mathrm{eff}},\mu)
=
f^{(y)}_{\mathrm{RT}}\!\left(y,z,U;\,\omega\rho,\;\omega\gamma,\;K_{\mathrm{eff}},\;\omega\mu\right).
\end{equation}
Hence, for any $c>0$, the pairs $(\omega,\theta)$ and
$\bigl(\omega/c,\;T_{c}\theta\bigr)$ with
$T_{c}:(\rho,\gamma,K_{\mathrm{eff}},\mu,\tau)\mapsto(c\rho,c\gamma,K_{\mathrm{eff}},c\mu,\tau)$
generate the \emph{identical} vector field, and $\{T_{c}\}_{c>0}$ is a
one-parameter group. Note that $K_{\mathrm{eff}}$ must \emph{not} be rescaled:
by Eq.~\eqref{eq:frt_expanded} scaling $\rho$ already scales the quadratic term,
and scaling $K_{\mathrm{eff}}$ too would undo it. The same applies to $\tau$,
which lives in the $\dot z$ equation and is untouched by $\omega$.

\paragraph{The residual weight.}
The closure's output layer is affine, $g_{\psi}=W_{3}h_{2}+b_{3}$, so for any
$k>0$
\begin{equation}
\label{eq:sr_gauge}
s_{r}\,g_{\psi}(\,\cdot\,;W_{1},b_{1},W_{2},b_{2},W_{3},b_{3})
=
\frac{s_{r}}{k}\,g_{\psi}(\,\cdot\,;W_{1},b_{1},W_{2},b_{2},kW_{3},kb_{3}).
\end{equation}

Both statements were verified numerically in double precision: the maximum
relative trajectory difference is $0$ for the $\omega$
transformation and $2.1\times10^{-16}$ for the $s_{r}$ transformation. With the
near-zero output initialisation prescribed in Section~\ref{sec:surrogates},
$s_{r}$ additionally has \emph{no} effect at all at initialisation: values
spanning five orders of magnitude give byte-identical trajectories.

\emph{Consequence.} $(\omega,s_{r})$ is a two-parameter gauge group acting on
$(\theta,\psi)$, not a mixing weight. The reason the two numbers cannot be
normalised to sum to one is not notational: both branches sweep out cones in
function space, so any rescaling of either weight is absorbed by that branch's
own parameters. An ablation over $(\omega,s_{r})$ therefore measures the
\emph{initialisation and optimisation path}, not the strength of the physics,
and the fitted $\rho,\gamma,\mu$ are gauge coordinates rather than physical
rates --- only $\omega\rho$, $\omega\gamma$, $\omega\mu$, $K_{\mathrm{eff}}$ and
$\tau$ are comparable between settings.

\subsection{A calibrated physics--closure control}
\label{sec:convex}

Making the weights convex is by itself insufficient: a cone remains a cone, so
$(1-\lambda)$ would simply be reabsorbed into $\theta$. We therefore combine a
convex mixture with two constraints that remove the two gauges:
\begin{equation}
\label{eq:cpinode}
\boxed{\;
\dot y \;=\; \underbrace{(1-\lambda)\,f^{(y)}_{\mathrm{RT}}(y,z,U;\theta)}_{\text{physics}}
\;+\;
\underbrace{\lambda\,\sigma_{\mathrm{ref}}\,\tanh\!\bigl(g_{\psi}(y,z,t,U)\bigr)}_{\text{closure}},
\qquad
\dot z = -\frac{z}{\tau}+U(t). \;}
\end{equation}
\begin{enumerate}
\item \textbf{Bounded closure.} $\sigma_{\mathrm{ref}}\tanh(\cdot)$ replaces the
  cone $\{s_{r}g_{\psi}\}$ by a ball of radius $\lambda\sigma_{\mathrm{ref}}$, so
  $\lambda$ genuinely caps the closure's authority and Eq.~\eqref{eq:sr_gauge}
  no longer applies. The scale $\sigma_{\mathrm{ref}}$ is estimated from the
  assimilation data alone, as the root-mean-square derivative of a
  Savitzky--Golay smooth of the sparse observations; measured against the true
  derivative RMS on the cohort it lands at $1.45\times$ of it, so the two
  branches of Eq.~\eqref{eq:cpinode} are commensurate and $\lambda$ means the
  same thing from patient to patient.
\item \textbf{Anchored physics.} A log-prior
  $w_{\theta}\bigl\|\log(\theta/\theta_{\mathrm{ref}})\bigr\|^{2}$, with
  $\theta_{\mathrm{ref}}$ the assimilated mechanistic fit, penalises exactly the
  rescaling $T_{c}$ of Eq.~\eqref{eq:omega_gauge} and so removes the remaining
  gauge. Without it the substitution $\lambda\to\lambda'$,
  $\theta\to\frac{1-\lambda}{1-\lambda'}\theta$ leaves the physics branch
  unchanged --- verified at $1.6\times10^{-16}$.
\item \textbf{Orthogonal closure.} A penalty
  $w_{\perp}\bigl\|P_{\mathrm{span}(\partial f/\partial\theta)}\tanh g_{\psi}\bigr\|^{2}$,
  with $P$ the least-squares projection onto the span of the mechanistic
  parameter sensitivities along the trajectory, prevents the closure from
  performing any correction that a change of $\theta$ could have performed. This
  is what identifies the physics--closure \emph{ratio}, and it directly targets
  the failure mode this work is about: reduced surrogates absorbing unresolved
  spatial effects into distorted effective parameters.
\end{enumerate}

Because $\lambda$ is now a bound on authority rather than a free scale, its two
endpoints are exact: $\lambda=0$ recovers the assimilated mechanistic model
identically, and $\lambda=1$ is a purely learned --- but bounded --- vector field
retaining the mechanistic damage state. The first identity is confirmed to the
last reported digit by the benchmark (Table~\ref{tab:main_results}).

\subsection{Learned gating}
\label{sec:gate}

A single scalar $\lambda$ forces the same physics--closure balance everywhere on
the trajectory, although the reduced physics is adequate in some regimes
(undisturbed growth) and inadequate in others (immediately after a fraction,
during recolonisation). We therefore also consider a state-dependent gate
\begin{equation}
\label{eq:gate}
\lambda(y,z,t,U) \;=\; \sigma\!\bigl(h_{\xi}(y,z,t,U) + \beta_{0}\bigr),
\qquad
\mathcal{L} \;\mathrel{+}=\; \eta\;\overline{\lambda},
\end{equation}
where $\sigma$ is the logistic function, $\beta_{0}$ sets the prior level, and
the penalty $\eta\overline{\lambda}$ on the trajectory-averaged gate expresses a
preference for the physics wherever the closure is not needed. Equations
\eqref{eq:cpinode} and \eqref{eq:gate} are otherwise identical.

\subsection{Reporting the realised balance}
\label{sec:phi}

Because a nominal weight need not equal the balance a fitted model actually
adopts, we report both. Along the integrated trajectory we measure the
\emph{realised physics share}
\begin{equation}
\label{eq:phi}
\Phi(t) \;=\;
\frac{\bigl|\dot y_{\mathrm{phys}}(t)\bigr|}
     {\bigl|\dot y_{\mathrm{phys}}(t)\bigr| + \bigl|\dot y_{\mathrm{ML}}(t)\bigr|}
\;\in\;[0,1],
\end{equation}
where $\dot y_{\mathrm{phys}}$ and $\dot y_{\mathrm{ML}}$ are the two additive
terms of Eq.~\eqref{eq:pinode_weighted} or Eq.~\eqref{eq:cpinode}. $\Phi$ is
defined identically for both parameterisations, which places them on a single
comparable axis, and it is reported separately inside the assimilation window
and in the forecast window.

One caveat is worth stating. $\Phi$ measures the two branches' magnitudes, not
their agreement: two large terms of opposite sign that nearly cancel give
$\Phi\approx\tfrac12$ while contributing little net drift. $\Phi$ is therefore an
honest measure of \emph{how much work each branch is doing}, which is the
quantity a mixing weight is supposed to control, but it should be read alongside
the forecast error rather than on its own.
